En un accident de trànsit es veu involucrat un camió que transporta diferents tipus de residus nuclears. A conseqüència de l'impacte, part de la càrrega ha perdut l'embalatge i no es pot identificar el material recollit. Per l'aspecte, s'ha pogut saber que només s'han vessat dos tipus diferents de residus, però no és possible identificar amb certesa quin és cadascun. Per tal de gestionar correctament el material recollit, els agents de neteja encarreguen un estudi de l'activitat radioactiva que tenen. Es prenen dues mostres d'1 µg i s'obtenen els resultats següents:

\begin{center}
\small
\begin{tabular}{lcccccccccc}
\toprule
Temps (min) & 60 & 150 & 240 & 330 & 420 & 510 & 600 & 690 & 780 & 870\\
Mostra A (Bq) & 1\,535 & 1\,025 & 678 & 484 & 315 & 238 & 139 & 103 & 78 & 45\\
Mostra B (Bq) & 260 & 215 & 170 & 145 & 114 & 92 & 82 & 68 & 55 & 47\\
\bottomrule
\end{tabular}
\end{center}

\begin{apartat}{1,25}
Representeu l'activitat en funció del temps. Feu un gràfic per a cada mostra en la quadrícula corresponent. A cada gràfic, dibuixeu-hi la línia de tendència de les dades.
\end{apartat}

\begin{apartat}{1,25}
Determineu, aproximadament, el període de semidesintegració de cada mostra.
\end{apartat}

Mostra A:
\figura[0.75]{fig1.png}

Mostra B:
\figura[0.75]{fig2.png}
